% Title Page
\makeisititle{EVALUASI KEPATUHAN HARMONI FUNGSIONAL OUTPUT MODEL GENERATIF MUSIK SIMBOLIK: STUDI KOMPARATIF TEORI GUSTAV STRUBE}{\researchername}{\researchernim}{Semester Ganjil 2026/2027}{2026}

% Approval Page
\newpage
\begin{spacing}{1.15}
\begin{center}
    HALAMAN PENGESAHAN\par
    \vspace{1cm}
    \textbf{EVALUASI KEPATUHAN HARMONI FUNGSIONAL OUTPUT MODEL GENERATIF MUSIK SIMBOLIK: STUDI KOMPARATIF TEORI GUSTAV STRUBE}\par
    \vspace{0.5cm}
    Diajukan Oleh:\par
    \researchername\par
    NIM: \researchernim\par
    \vspace{1.5cm}
    Proposal ini telah disetujui sebagai syarat melaksanakan Tugas Akhir/Skripsi di\par
    \programstudy, \facultyname,\par
    \institutionname\par
    \vspace{2cm}
    
    \begin{minipage}[t]{0.45\textwidth}
        \centering
        ~\\
        Mengetahui,\par
        \vspace{3cm}
        \textbf{\advisoracademic}\par
        NIP. \advisoracademicnip
    \end{minipage}
    \hfill
    \begin{minipage}[t]{0.45\textwidth}
        \centering
        Yogyakarta, \makebox[3cm]{\hrulefill}\\
        Disetujui,\par
        \vspace{3cm}
        \textbf{\advisorthesis}\par
        NIP. \advisorthesisnip
    \end{minipage}
\end{center}
\end{spacing}

% Abstract
\newpage
\begin{spacing}{1.15}
\begin{center}
    \textbf{ABSTRAK}
\end{center}

Penelitian ini mengevaluasi tingkat kepatuhan model kecerdasan buatan (AI) generatif musik simbolik terhadap kaidah harmoni fungsional Barat berdasarkan teori Gustav Strube. Tiga model dari era berbeda dievaluasi: DeepBach (2017), Coconet (2017), dan NotaGen (2025). Evaluasi dilakukan secara algoritmik menggunakan pustaka \textit{music21} dengan kriteria kuint sejajar, oktaf sejajar, dan resolusi nada penuntun. Hasil penelitian diharapkan dapat menunjukkan sejauh mana perkembangan arsitektur AI musik mampu memahami aturan deterministik dalam harmoni fungsional.\par

\textbf{Kata Kunci:} Musik Simbolik, AI, Harmoni Fungsional, Gustav Strube, Evaluasi.
\end{spacing}

% Abstract (English)
\newpage
\begin{spacing}{1.15}
\begin{center}
    \textbf{ABSTRACT}
\end{center}

This study evaluates the compliance rate of symbolic music generative Artificial Intelligence (AI) models with Western functional harmony rules based on Gustav Strube's theory. Three models from different eras are evaluated: DeepBach (2017), Coconet (2017), and NotaGen (2025). The evaluation is performed algorithmically using the \textit{music21} library with criteria of parallel fifths, parallel octaves, and leading-tone resolution. The results of the study are expected to show the extent to which the development of music AI architectures is capable of understanding deterministic rules in functional harmony.\par

\textbf{Keywords:} Symbolic Music, AI, Functional Harmony, Gustav Strube, Evaluation.
\end{spacing}
\newpage
